\section{Introduction}\label{sec:intro}

\begin{figure}[t]
\centering
\resizebox{\linewidth}{!}{%
\begin{tikzpicture}[
  font=\footnotesize,
  box/.style={draw, rounded corners=2pt, align=center, minimum height=9mm, inner sep=3pt, fill=white},
  eng/.style={box, fill=blue!7, draw=blue!60!black, thick},
  geo/.style={box, fill=orange!10, draw=orange!70!black},
  ros/.style={box, fill=gray!8, draw=gray!70!black},
  arr/.style={-{Latex[length=2mm]}, thick},
  lab/.style={font=\scriptsize\itshape, text=black!70}]
% ---- engine (geometry-free) ----
\node[eng, minimum width=10.4cm, minimum height=15mm] (lm) at (0,0)
  {\textbf{\code{LayeredMap} + \code{PeriodicityModel}} \quad (shared, keyed only by \code{int64} \code{CellId})\\[2pt]
   window counters $\to$ log-odds $+$ survival decay $\lambda$ $\to$ Schmitt trigger $[\pdem,\pgrad]$ $\to$ FreMEn-lite amplitude $a$ $\to$ prune\\[1pt]
   \emph{Static} \;/\; \emph{Periodic} \;/\; \emph{Transient} \;/\; \emph{Unknown}};
% ---- backends ----
\node[geo, below left=11mm and -48mm of lm, minimum width=4.9cm] (g2)
  {\textbf{\code{Grid2DBackend}}\\ $\kappa_{2\mathrm{D}}$: row-major array index\\ ray walk: integer Bresenham};
\node[geo, below right=11mm and -48mm of lm, minimum width=4.9cm] (v3)
  {\textbf{\code{Voxel3DBackend}}\\ $\kappa_{3\mathrm{D}}$: 3$\times$21-bit packed hash key\\ ray walk: half-voxel sample march};
\draw[arr] (g2.north) -- node[lab, left] {\code{observeHit/Miss}(id)} (g2.north |- lm.south);
\draw[arr] (v3.north) -- node[lab, right] {\code{observeHit/Miss}(id)} (v3.north |- lm.south);
% band labels
\begin{scope}[on background layer]
\node[fit=(lm)(g2)(v3), draw=blue!40, dashed, rounded corners=4pt, inner sep=6pt,
  label={[lab, anchor=north east]north east:\code{strata\_core}: C++17 + Eigen, no ROS/PCL, 41 gtest cases}] (core) {};
\end{scope}
% ---- ROS node ----
\node[ros, below=12mm of core.south, minimum width=10.4cm] (node)
  {\textbf{\code{strata::MappingNode}} (ROS~2 Humble)\quad \code{backend:=grid2d|voxel3d} selects one \code{unique\_ptr} at construction\\
   \code{LaserScan} / \code{PointCloud2} $\to$ \code{Observation} (map frame) \qquad
   publishes \code{\textasciitilde/map} or \code{\textasciitilde/map\_points}, serves \code{\textasciitilde/save\_map}};
\draw[arr] (node.north -| g2.south) -- node[lab, left] {\code{integrate}, \code{tick}} (g2.south);
\draw[arr] (node.north -| v3.south) -- node[lab, right] {\code{integrate}, \code{tick}} (v3.south);
\node[ros, below=6mm of node, minimum width=4.2cm, fill=white, dashed] (tf) {external localizer (TF: \code{map} $\to$ sensor)};
\draw[arr, dashed] (tf) -- (node);
\end{tikzpicture}}
\caption{\textbf{Overview of \sname.} One geometry-free engine (top) owns every persistence and periodicity decision and sees only integer cell keys. The two backends (middle) contribute exactly two things: the point-to-key map $\kappa_b$ and the free-space ray walk. A thin ROS~2 node (bottom) selects one backend from a single string parameter and consumes, but never produces, the sensor pose.}
\label{fig:overview}
\end{figure}


An occupancy grid built the textbook way accumulates evidence and never lets go of it \citep{thrun2005probabilistic}. For a robot that drives a route once this is harmless; for one that maps the same warehouse or hospital corridor for weeks it is a slow failure. A cart parked for an afternoon, a person walking past and a load-bearing wall all leave the same kind of mark, and after enough passes the map describes the union of everything the sensor ever touched rather than the building. Lifelong mapping has to hold three requirements at once: the map must be \emph{durable} enough to keep structure that is genuinely permanent, \emph{plastic} enough to forget structure that has moved on, and it must do both without one global forgetting rate that is simultaneously too slow to erase movers and too fast to trust a wall---the stability--plasticity dilemma that \citet{biber2005dynamicmaps} framed for dynamic maps. A third requirement complicates the split: some structure is \emph{cyclic}, such as a door open every morning and shut every night. Its occupancy is a periodic temporal signal, and a flat decay law erases exactly the regularity that makes it predictable \citep{krajnik2017fremen,krajnik2014spectral}.

The systems that come closest to meeting all three are heavier than the problem requires. ELite \citep{gil2025elite} and LT-mapper \citep{kim2022ltmapper} resolve durability against plasticity, but as stages of full lifelong-SLAM pipelines with multi-session registration, place recognition and alignment. RTAB-Map \citep{labbe2019rtabmap} is the closest engineering precedent for selectable 2D and 3D geometry under one framework, and it too is a complete SLAM system with loop closure and memory tiering. Khronos \citep{schmid2024khronos} generalises the two-timescale idea to a metric-semantic scene graph. In each case the persistence-and-classification logic a robot needs is welded inside a larger apparatus and, usually, to one map geometry. What the surveyed prior art lacks is that piece on its own: a standalone, SLAM-free, dependency-light engine that performs only graduation, forgetting and periodicity labelling, runs under either 2D or 3D geometry, and can sit beneath an external localizer.

To overcome this coupling, we introduce \emph{\sname}, a persistence-and-periodicity engine whose cells carry temporal strata---Static, Periodic and Transient classes over an Unknown baseline---inside one per-cell state machine (Figure~\ref{fig:overview}). The state machine never sees coordinates. Both shipped backends translate world points into \code{int64} cell keys and walk a free-space ray to clear it; everything downstream of the key happens once, in the shared engine. Choosing 2D or 3D is one runtime string parameter resolved at construction, and the engine (\code{strata\_core}) links no ROS, DDS, TF or PCL, so its scientific logic builds and unit-tests without a ROS installation. We state the scope up front: \sname\ is not SLAM (it consumes an external \code{map}$\to$sensor transform), it models single-session persistence rather than multi-session remapping, and its evaluation is synthetic. Our contributions are:
\begin{itemize}
  \item \emph{One Geometry-Free Engine, Two Geometries}: a single \code{int64}-keyed engine (\code{LayeredMap} with \code{PeriodicityModel}) drives a fixed-array 2D grid and an unbounded 3D voxel hash through one \code{MapBackend} interface; the only per-backend code is the point-to-key map and the free-space ray walk (\S\ref{sec:method}), and the core builds with C++17 and Eigen alone under 41 behaviour-level tests.
  \item \emph{A Minimal Lifelong Classifier}: survival decay after the Persistence Filter \citep{rosen2016persistence}, a Schmitt-trigger graduation band motivated by Removert \citep{kim2020removert}, negative-evidence ray clearing as in ReFusion \citep{palazzolo2019refusion}, and a parallel FreMEn-lite periodicity test \citep{krajnik2017fremen} whose per-read-out false-alarm probability under iid clutter and occupancy-independent sampling is bounded for any noise rate by a Chernoff argument (Proposition~\ref{prop:chernoff}), emitting a four-class label from one SLAM-free module, with every equation transcribed from the shipped code and three implementation-versus-specification differences stated explicitly.
  \item \emph{A Mechanistic, Reproducible Characterisation}: a calibration run on held-out seeds (E0) and four seeded experiments (E1--E4) that isolate backend equivalence, periodicity detection, the contribution of hysteresis and decay, and per-dimension cost, with every table and figure regenerated from committed CSVs and failure modes diagnosed rather than averaged away.
\end{itemize}

Across the suite, both backends graduate a 161-cell wall by the third window and keep recall at 1.0 through window 39; their only quality gap, a final static precision of 0.9253 (\code{grid2d}) against 0.9758 (\code{voxel3d}) under 100 movers per window, arises in the backends' native sampling rather than in the shared classifier. We report and correct two defects of the first release's periodicity test, with before-and-after results. An uncentred Fourier amplitude could demote constant walls. An uncalibrated amplitude threshold labelled Bernoulli clutter Periodic at 0 to 4 of 5 cells depending on read-out length, a mean false-positive rate of 0.331 even after centring. With the calibrated test, the false-positive rate is 0 at every read-out length from 8 to 100 windows. Both 50\%-duty doors are still detected, for a true-positive rate of 2/3, but only from 15 windows instead of 8, and the one miss is a provable interaction between pruning and the amplitude gate. Removing the hysteresis band raises flicker from 54 to 574 toggles at the same decay, and the engine costs a flat 56~B per cell, with the 5.5--15.0$\times$ per-call gap between backends explained by 4.4--10.0$\times$ more live 3D cells. These results show that the lifelong-mapping core can be written once and shared across dimensions without a per-dimension cost of its own.
